\section{Diagonal club intersections and Fodor's lemma}

\import{stationary/club.tex}
\import{set/function-partition.tex}

The indexed intersection theorem below exposes both bounds that keep the
usual hand-over-hand construction inside its ambient ordinal: the family has
length below the cofinality, and the final iteration is countable.

\begin{definition}\label{club_family_choice}
    $\clubFamilyChoice{F}{I}{\beta}{\alpha}
        =\{(i,\choice{\clubUpper{F(i)}{\beta}{\alpha}})\mid i\in I\}$.
\end{definition}

\begin{definition}\label{club_family_advance}
    $\clubFamilyAdvance{F}{I}{\alpha}
        =\{(\beta,\suc{\unions{\ran{
            \clubFamilyChoice{F}{I}{\beta}{\alpha}}}})
            \mid \beta\in\alpha\}$.
\end{definition}

\begin{definition}\label{diagonal_club_advance}
    $\diagonalClubAdvance{F}{\alpha}
        =\{(\beta,
            \choice{\clubUpper{
                \indexedIntersection{\suc{\beta}}{F}{\alpha}}
                {\beta}{\alpha}})
            \mid \beta\in\alpha\}$.
\end{definition}

\begin{definition}\label{fodor_avoiding_clubs}
    $\fodorAvoidingClubs{f}{\eta}{\kappa}
        =\{C\in\pow{\kappa}\mid
            \clubIn{C}{\kappa}\land
            \fibre{f}{\eta}\notmeets C\}$.
\end{definition}

\begin{definition}\label{fodor_club_family}
    $\fodorClubFamily{f}{\kappa}
        =\{(\eta,\choice{\fodorAvoidingClubs{f}{\eta}{\kappa}})
            \mid \eta\in\kappa\}$.
\end{definition}

\begin{definition}\label{fodor_values}
    $\fodorValues{f}{\kappa}
        =\{\eta\in\kappa\mid
            \text{$\fibre{f}{\eta}$ is stationary in $\kappa$}\}$.
\end{definition}

\begin{definition}\label{fodor_value}
    $\fodorValue{f}{\kappa}=\inters{\fodorValues{f}{\kappa}}$.
\end{definition}

\begin{definition}\label{fodor_fibre}
    $\fodorFibre{f}{\kappa}=\fibre{f}{\fodorValue{f}{\kappa}}$.
\end{definition}

\begin{definition}\label{regressive_on}
    $f$ is regressive on $S$ iff $f$ is a function on $S$ and
        for every $\gamma\in S\setminus\{\emptyset\}$ we have
            $f(\gamma)\in\gamma$.
\end{definition}

\begin{proposition}\label{ordinal_closed_indexed_intersection}
    Let $\alpha$ be an ordinal.
    Suppose for every $i\in I$ we have
        $\ordinalClosed{F(i)}{\alpha}$.
    Then $\ordinalClosed{\indexedIntersection{I}{F}{\alpha}}{\alpha}$.
\end{proposition}
\begin{proof}
    Let $C=\indexedIntersection{I}{F}{\alpha}$.
    $C\subseteq\alpha$ by \auto{indexed_intersection_subseteq}.
    Show that for every limit ordinal $\gamma\in\alpha$ such that
        $\ordinalUnbounded{C\inter\gamma}{\gamma}$
        we have $\gamma\in C$.
    \begin{subproof}
        Fix $\gamma$.
        Assume $\gamma\in\alpha$.
        Assume $\gamma$ is a limit ordinal.
        Assume $\ordinalUnbounded{C\inter\gamma}{\gamma}$.
        Show that for every $i\in I$ we have $\gamma\in F(i)$.
        \begin{subproof}
            Fix $i\in I$.
            $C\subseteq F(i)$
                by \auto{indexed_intersection_member_subseteq}.
            Hence $C\inter\gamma\subseteq F(i)\inter\gamma$
                by \obv{subseteq,inter}.
            $F(i)\inter\gamma\subseteq\gamma$
                by \obv{subseteq,inter}.
            Hence $\ordinalUnbounded{F(i)\inter\gamma}{\gamma}$
                by \obv{ordinal_unbounded_superset}.
            Follows by \auto{ordinal_closed}.
        \end{subproof}
        Follows by \obv{indexed_intersection}.
    \end{subproof}
    Follows by \auto{ordinal_closed}.
\end{proof}

\begin{proposition}\label{club_family_choice_type}
    Let $\alpha$ be a limit ordinal.
    Suppose $F$ is a function and $I\subseteq\dom{F}$.
    Suppose for every $i\in I$ we have $\clubIn{F(i)}{\alpha}$.
    Let $\beta\in\alpha$.
    Let $q=\clubFamilyChoice{F}{I}{\beta}{\alpha}$.
    Then $q\in\funs{I}{\alpha}$ and for every $i\in I$ we have
        $q(i)\in F(i)$ and $\beta\in q(i)$.
\end{proposition}
\begin{proof}
    $q$ is a relation by \auto{relation,club_family_choice}.
    $q$ is right-unique
        by \auto{rightunique,club_family_choice,pair_eq_iff}.
    $\dom{q}=I$
        by \auto{setext,dom_iff,club_family_choice,pair_eq_iff}.
    Show that for every $i\in I$ we have
        $q(i)\in F(i)$ and $\beta\in q(i)$ and $q(i)\in\alpha$.
    \begin{subproof}
        Fix $i\in I$.
        Hence $i\in\dom{F}$ by \obv{subseteq}.
        $\ordinalUnbounded{F(i)}{\alpha}$
            by \auto{club_in_elim}.
        Hence $\clubUpper{F(i)}{\beta}{\alpha}$ is inhabited
            by \obv{club_upper_inhabited}.
        Let $x=\choice{\clubUpper{F(i)}{\beta}{\alpha}}$.
        $x\in\clubUpper{F(i)}{\beta}{\alpha}$ by \auto{choice_in}.
        Hence $x\in F(i)$ and $\beta\in x$
            by \obv{club_upper}.
        $F(i)\subseteq\alpha$ by \obv{club_in_elim}.
        Hence $x\in\alpha$ by \obv{subseteq}.
        $(i,x)\in q$ by \obv{club_family_choice}.
        Hence $q(i)=x$ by \obv{function_apply_intro}.
    \end{subproof}
    Hence $q$ is a function to $\alpha$ by \auto.
    Hence $q\in\funs{I}{\alpha}$ by \obv{funs_intro}.
\end{proof}

\begin{theorem}\label{club_family_advance_type}
    Let $\alpha$ be a limit ordinal.
    Let $I\in\cofinality{\alpha}$.
    Suppose $F$ is a function and $I\subseteq\dom{F}$.
    Suppose for every $i\in I$ we have $\clubIn{F(i)}{\alpha}$.
    Then $\clubFamilyAdvance{F}{I}{\alpha}
        \in\funs{\alpha}{\alpha}$.
\end{theorem}
\begin{proof}
    Let $h=\clubFamilyAdvance{F}{I}{\alpha}$.
    $h$ is a relation by \auto{relation,club_family_advance}.
    $h$ is right-unique
        by \auto{rightunique,club_family_advance,pair_eq_iff}.
    $\dom{h}=\alpha$
        by \auto{setext,dom_iff,club_family_advance,pair_eq_iff}.
    Show that for every $\beta\in\dom{h}$ we have $h(\beta)\in\alpha$.
    \begin{subproof}
        Fix $\beta\in\dom{h}$.
        Hence $\beta\in\alpha$.
        Let $q=\clubFamilyChoice{F}{I}{\beta}{\alpha}$.
        $q\in\funs{I}{\alpha}$ by \auto{club_family_choice_type}.
        Let $U=\unions{\ran{q}}$.
        $U\in\alpha$ by \auto{small_domain_range_union_below}.
        Hence $\suc{U}\in\alpha$ by \obv{suc_elem_limitordinal}.
        $(\beta,\suc{U})\in h$ by \auto{club_family_advance}.
        Hence $h(\beta)=\suc{U}$ by \obv{function_apply_intro}.
    \end{subproof}
    Hence $h$ is a function to $\alpha$.
    Follows by \auto{funs_intro}.
\end{proof}

\begin{proposition}\label{club_family_advance_properties}
    Let $\alpha$ be a limit ordinal.
    Let $I\in\cofinality{\alpha}$.
    Suppose $I$ is inhabited.
    Suppose $F$ is a function and $I\subseteq\dom{F}$.
    Suppose for every $i\in I$ we have $\clubIn{F(i)}{\alpha}$.
    Let $\beta\in\alpha$.
    Let $q=\clubFamilyChoice{F}{I}{\beta}{\alpha}$.
    Let $U=\unions{\ran{q}}$.
    Let $h=\clubFamilyAdvance{F}{I}{\alpha}$.
    Then $h(\beta)=\suc{U}$ and $\beta\in h(\beta)$ and
        for every $i\in I$ we have
            $q(i)\in F(i)$ and $\beta\in q(i)$ and
            $q(i)\in h(\beta)$.
\end{proposition}
\begin{proof}
    $h\in\funs{\alpha}{\alpha}$
        by \auto{club_family_advance_type}.
    Hence $h$ is a function by \obv{funs_is_function}.
    $(\beta,\suc{U})\in h$ by \auto{club_family_advance}.
    Hence $h(\beta)=\suc{U}$ by \obv{function_apply_intro}.
    $q\in\funs{I}{\alpha}$ and for every $i\in I$ we have
        $q(i)\in F(i)$ and $\beta\in q(i)$
        by \obv{club_family_choice_type}.
    Hence $q$ is a function and $\dom{q}=I$
        by \obv{funs_is_function,funs}.
    $U\in\alpha$ by \auto{small_domain_range_union_below}.
    Hence $U$ is an ordinal
        by \obv{limit_ordinal,elem_of_ordinal_is_ordinal}.
    Show that for every $i\in I$ we have $q(i)\in h(\beta)$.
    \begin{subproof}
        Fix $i\in I$.
        $q(i)\in\ran{q}$ by \auto{fun_ran_iff}.
        Hence $q(i)\subseteq U$ by \auto{subseteq,unions_iff}.
        $q(i)\in\alpha$ by \obv{funs_type_apply}.
        Hence $q(i)$ is an ordinal
            by \obv{limit_ordinal,elem_of_ordinal_is_ordinal}.
        Hence $q(i)\in U$ or $q(i)=U$
            by \obv{ordinal_subseteq_implies_elem_or_eq}.
        Hence $q(i)\in\suc{U}$ by \auto{suc_iff}.
    \end{subproof}
    Take $i_0$ such that $i_0\in I$.
    $\beta\in q(i_0)$ by \auto{club_family_choice_type}.
    $q(i_0)\in h(\beta)$.
    $h(\beta)\in\alpha$ by \obv{funs_type_apply}.
    Hence $h(\beta)$ is an ordinal
        by \obv{limit_ordinal,elem_of_ordinal_is_ordinal}.
    Hence $\beta\in h(\beta)$ by \obv{ordinal_transitivity}.
\end{proof}

\begin{proposition}\label{club_family_iteration_properties}
    Let $\alpha$ be a limit ordinal.
    Let $I\in\cofinality{\alpha}$.
    Suppose $I$ is inhabited.
    Suppose $F$ is a function and $I\subseteq\dom{F}$.
    Suppose for every $i\in I$ we have $\clubIn{F(i)}{\alpha}$.
    Let $\beta\in\alpha$.
    Let $h=\clubFamilyAdvance{F}{I}{\alpha}$.
    Let $a=h(\beta)$.
    Let $s=\iterationSet{h}{a}$.
    Then $s\in\funs{\naturals}{\alpha}$ and
        for every natural number $n$ we have
            $s(n)\in s(\suc{n})$ and for every $i\in I$ we have
                $(\clubFamilyChoice{F}{I}{s(n)}{\alpha})(i)\in F(i)$ and
                $s(n)\in(\clubFamilyChoice{F}{I}{s(n)}{\alpha})(i)$ and
                $(\clubFamilyChoice{F}{I}{s(n)}{\alpha})(i)
                    \in s(\suc{n})$.
\end{proposition}
\begin{proof}
    $h\in\funs{\alpha}{\alpha}$
        by \auto{club_family_advance_type}.
    $a\in\alpha$ by \auto{funs_type_apply}.
    $s\in\funs{\naturals}{\alpha}$
        by \auto{iteration_set_type}.
    Hence $s$ is a function and $\dom{s}=\naturals$
        by \obv{funs_is_function,funs}.
    Show that for every natural number $n$ we have
        $s(n)\in s(\suc{n})$ and for every $i\in I$ we have
            $(\clubFamilyChoice{F}{I}{s(n)}{\alpha})(i)\in F(i)$ and
            $s(n)\in(\clubFamilyChoice{F}{I}{s(n)}{\alpha})(i)$ and
            $(\clubFamilyChoice{F}{I}{s(n)}{\alpha})(i)
                \in s(\suc{n})$.
    \begin{subproof}
        Fix $n$.
        Assume $n$ is a natural number.
        $s(n)\in\alpha$ by \obv{funs_type_apply}.
        $s(\suc{n})=h(s(n))$
            by \auto{iteration_set_successor}.
        Let $q=\clubFamilyChoice{F}{I}{s(n)}{\alpha}$.
        $s(n)\in h(s(n))$ and for every $i\in I$ we have
            $q(i)\in F(i)$ and $s(n)\in q(i)$ and
            $q(i)\in h(s(n))$
            by \auto{club_family_advance_properties}.
        Follows by \auto.
    \end{subproof}
\end{proof}

\begin{theorem}\label{indexed_clubs_have_common_limit_above}
    Let $\alpha$ be a limit ordinal.
    Suppose $\naturals\in\cofinality{\alpha}$.
    Let $I\in\cofinality{\alpha}$.
    Suppose $I$ is inhabited.
    Suppose $F$ is a function and $I\subseteq\dom{F}$.
    Suppose for every $i\in I$ we have $\clubIn{F(i)}{\alpha}$.
    Let $\beta\in\alpha$.
    Then there exists
        $x\in\indexedIntersection{I}{F}{\alpha}$ such that
        $\beta\in x$ and $x$ is a limit ordinal.
\end{theorem}
\begin{proof}
    $\alpha$ is an ordinal by \obv{limit_ordinal}.
    Let $h=\clubFamilyAdvance{F}{I}{\alpha}$.
    $h\in\funs{\alpha}{\alpha}$
        by \auto{club_family_advance_type}.
    Let $a=h(\beta)$.
    $a\in\alpha$ by \auto{funs_type_apply}.
    Let $s=\iterationSet{h}{a}$.
    $s\in\funs{\naturals}{\alpha}$ and
        for every natural number $n$ we have
            $s(n)\in s(\suc{n})$ and for every $i\in I$ we have
                $(\clubFamilyChoice{F}{I}{s(n)}{\alpha})(i)\in F(i)$ and
                $s(n)\in(\clubFamilyChoice{F}{I}{s(n)}{\alpha})(i)$ and
                $(\clubFamilyChoice{F}{I}{s(n)}{\alpha})(i)
                    \in s(\suc{n})$
        by \obv{club_family_iteration_properties}.
    Hence $s$ is a function and $\dom{s}=\naturals$
        by \obv{funs_is_function,funs}.
    Let $x=\unions{\ran{s}}$.
    $x\in\alpha$ by \auto{small_domain_range_union_below}.
    $\ran{s}\subseteq\alpha$ by \obv{funs_ran}.
    Hence every element of $\ran{s}$ is an ordinal
        by \obv{subseteq,elem_of_ordinal_is_ordinal}.
    Hence $x$ is an ordinal by \obv{union_of_ordinals_is_ordinal}.
    $\emptyset$ is a natural number by \obv{omega_empty}.
    Hence $s(\emptyset)\in s(\suc{\emptyset})$
        by \auto{club_family_iteration_properties}.
    $\suc{\emptyset}$ is a natural number by \obv{omega_successor}.
    Hence $s(\suc{\emptyset})\in\ran{s}$ by \auto{fun_ran_iff}.
    Hence $s(\emptyset)\in x$ by \auto{unions_iff}.
    Hence $x\neq\emptyset$ by \auto{emptyset}.
    Hence $\emptyset\in x$
        by \obv{ordinal_empty_or_emptyset_elem}.
    Show that for every $\rho\in x$ we have $\suc{\rho}\in x$.
    \begin{subproof}
        Fix $\rho\in x$.
        Take $y\in\ran{s}$ such that $\rho\in y$
            by \obv{unions_iff}.
        Take $n\in\dom{s}$ such that $s(n)=y$
            by \obv{fun_ran_iff}.
        Hence $n$ is a natural number.
        Hence $s(n)\in s(\suc{n})$
            by \auto{club_family_iteration_properties}.
        $s(n)$ is an ordinal
            by \obv{funs_type_apply,elem_of_ordinal_is_ordinal}.
        Hence $\rho$ is an ordinal by \obv{elem_of_ordinal_is_ordinal}.
        Hence $\suc{\rho}\subseteq s(n)$
            by \obv{ordinal_suc_subseteq}.
        $\suc{\rho}$ is an ordinal by \obv{suc_ordinal}.
        Hence $\suc{\rho}\in s(n)$ or $\suc{\rho}=s(n)$
            by \obv{ordinal_subseteq_implies_elem_or_eq}.
        $s(\suc{n})$ is an ordinal
            by \obv{funs_type_apply,omega_successor,elem_of_ordinal_is_ordinal}.
        Hence $\suc{\rho}\in s(\suc{n})$
            by \auto{ordinal_transitivity}.
        Hence $\suc{n}\in\dom{s}$ by \obv{funs,omega_successor}.
        Hence $s(\suc{n})\in\ran{s}$ by \auto{fun_ran_iff}.
        Follows by \auto{unions_iff}.
    \end{subproof}
    Hence $\emptyset\precedes x$.
    Hence $x$ is a limit ordinal by \obv{limit_ordinal}.
    Show that for every $i\in I$ we have $x\in F(i)$.
    \begin{subproof}
        Fix $i\in I$.
        $\ordinalClosed{F(i)}{\alpha}$ and $F(i)\subseteq\alpha$
            by \auto{club_in_elim}.
        Show that $\ordinalUnbounded{F(i)\inter x}{x}$.
        \begin{subproof}
            $F(i)\inter x\subseteq x$ by \obv{subseteq,inter}.
            Show that for every $\rho\in x$ there exists
                $z\in F(i)\inter x$ such that $\rho\in z$.
            \begin{subproof}
                Fix $\rho\in x$.
                Take $y\in\ran{s}$ such that $\rho\in y$
                    by \obv{unions_iff}.
                Take $n\in\dom{s}$ such that $s(n)=y$
                    by \obv{fun_ran_iff}.
                Hence $n$ is a natural number.
                Let $q=\clubFamilyChoice{F}{I}{s(n)}{\alpha}$.
                $q(i)\in F(i)$ and $s(n)\in q(i)$ and
                    $q(i)\in s(\suc{n})$
                    by \auto{club_family_iteration_properties}.
                $\suc{n}$ is a natural number
                    by \obv{omega_successor}.
                Hence $s(\suc{n})\in\ran{s}$
                    by \auto{fun_ran_iff}.
                Hence $q(i)\in x$ by \auto{unions_iff}.
                Hence $q(i)\in F(i)\inter x$ by \obv{inter_intro}.
                $q(i)\in\alpha$ by \obv{subseteq}.
                Hence $q(i)$ is an ordinal
                    by \obv{elem_of_ordinal_is_ordinal}.
                $\rho\in s(n)$.
                Hence $\rho\in q(i)$ by \obv{ordinal_transitivity}.
            \end{subproof}
            Follows by \obv{ordinal_unbounded}.
        \end{subproof}
        Follows by \obv{ordinal_closed}.
    \end{subproof}
    Hence $x\in\indexedIntersection{I}{F}{\alpha}$
        by \obv{indexed_intersection}.
    $\beta\in a$ by \auto{club_family_advance_properties}.
    $s(\emptyset)=a$ by \auto{iteration_set_zero}.
    Hence $\beta\in s(\emptyset)$.
    Hence $\beta\in x$ by \obv{ordinal_transitivity}.
\end{proof}

\begin{theorem}\label{indexed_club_intersection}
    Let $\alpha$ be a limit ordinal.
    Suppose $\naturals\in\cofinality{\alpha}$.
    Let $I\in\cofinality{\alpha}$.
    Suppose $I$ is inhabited.
    Suppose $F$ is a function and $I\subseteq\dom{F}$.
    Suppose for every $i\in I$ we have $\clubIn{F(i)}{\alpha}$.
    Then $\clubIn{\indexedIntersection{I}{F}{\alpha}}{\alpha}$.
\end{theorem}
\begin{proof}
    Let $C=\indexedIntersection{I}{F}{\alpha}$.
    $C\subseteq\alpha$ by \auto{indexed_intersection_subseteq}.
    Show that $\ordinalUnbounded{C}{\alpha}$.
    \begin{subproof}
        Show that for every $\beta\in\alpha$ there exists $x\in C$
            such that $\beta\in x$.
        \begin{subproof}
            Fix $\beta\in\alpha$.
            Take $x\in C$ such that $\beta\in x$ and
                $x$ is a limit ordinal
                by \auto{indexed_clubs_have_common_limit_above}.
            Follows by \auto.
        \end{subproof}
        Follows by \obv{ordinal_unbounded}.
    \end{subproof}
    For every $i\in I$ we have $\ordinalClosed{F(i)}{\alpha}$
        by \auto{club_in_elim}.
    $\ordinalClosed{\indexedIntersection{I}{F}{\alpha}}{\alpha}$
        by \obv{ordinal_closed_indexed_intersection}.
    Hence $\ordinalClosed{C}{\alpha}$.
    Follows by \auto{club_in}.
\end{proof}

\begin{proposition}\label{diagonal_intersection_closed}
    Let $\alpha$ be an ordinal.
    Suppose for every $i\in\alpha$ we have
        $\ordinalClosed{F(i)}{\alpha}$.
    Then $\ordinalClosed{\diagonalIntersection{F}{\alpha}}{\alpha}$.
\end{proposition}
\begin{proof}
    Let $D=\diagonalIntersection{F}{\alpha}$.
    $D\subseteq\alpha$ by \auto{diagonal_intersection,subseteq}.
    Show that for every limit ordinal $\gamma\in\alpha$ such that
        $\ordinalUnbounded{D\inter\gamma}{\gamma}$
        we have $\gamma\in D$.
    \begin{subproof}
        Fix $\gamma$.
        Assume $\gamma\in\alpha$.
        Assume $\gamma$ is a limit ordinal.
        Assume $\ordinalUnbounded{D\inter\gamma}{\gamma}$.
        Show that for every $i\in\gamma$ we have $\gamma\in F(i)$.
        \begin{subproof}
            Fix $i\in\gamma$.
            Hence $i\in\alpha$ by \auto{ordinal_transitivity}.
            $\ordinalClosed{F(i)}{\alpha}$.
            Show that $\ordinalUnbounded{F(i)\inter\gamma}{\gamma}$.
            \begin{subproof}
                $F(i)\inter\gamma\subseteq\gamma$
                    by \obv{subseteq,inter}.
                Show that for every $\rho\in\gamma$ there exists
                    $\delta\in F(i)\inter\gamma$ such that
                    $\rho\in\delta$.
                \begin{subproof}
                    Fix $\rho\in\gamma$.
                    $i$ and $\rho$ are ordinals
                        by \obv{limit_ordinal,elem_of_ordinal_is_ordinal}.
                    $\rho\union i$ is an ordinal
                        by \obv{union_of_two_ordinals_is_ordinal}.
                    $\rho\union i\in\gamma$ by \obv{ordinal_union_below}.
                    Take $\delta\in D\inter\gamma$ such that
                        $\rho\union i\in\delta$ by \obv{ordinal_unbounded}.
                    Hence $\delta\in D$ and $\delta\in\gamma$
                        by \obv{inter_elim_left,inter_elim_right}.
                    $\delta$ is an ordinal
                        by \obv{limit_ordinal,elem_of_ordinal_is_ordinal}.
                    $i\subseteq\rho\union i$ and $\rho\subseteq\rho\union i$
                        by \auto{union_upper_left,union_upper_right}.
                    Hence $i\in\rho\union i$ or $i=\rho\union i$
                        by \auto{ordinal_subseteq_implies_elem_or_eq}.
                    Hence $i\in\delta$ by \auto{ordinal_transitivity}.
                    Hence $\delta\in F(i)$
                        by \auto{diagonal_intersection}.
                    Hence $\delta\in F(i)\inter\gamma$
                        by \obv{inter_intro}.
                    Hence $\rho\in\rho\union i$ or $\rho=\rho\union i$
                        by \auto{ordinal_subseteq_implies_elem_or_eq}.
                    Hence $\rho\in\delta$ by \auto{ordinal_transitivity}.
                    Follows by \auto.
                \end{subproof}
                Follows by \obv{ordinal_unbounded}.
            \end{subproof}
            Follows by \obv{ordinal_closed}.
        \end{subproof}
        Follows by \obv{diagonal_intersection}.
    \end{subproof}
    Follows by \auto{ordinal_closed}.
\end{proof}

\begin{theorem}\label{diagonal_club_advance_type}
    Let $\alpha$ be an ordinal.
    Suppose for every $\beta\in\alpha$ we have
        $\ordinalUnbounded{
            \indexedIntersection{\suc{\beta}}{F}{\alpha}}{\alpha}$.
    Then $\diagonalClubAdvance{F}{\alpha}\in\funs{\alpha}{\alpha}$.
\end{theorem}
\begin{proof}
    Let $h=\diagonalClubAdvance{F}{\alpha}$.
    $h$ is a relation by \auto{relation,diagonal_club_advance}.
    $h$ is right-unique
        by \auto{rightunique,diagonal_club_advance,pair_eq_iff}.
    $\dom{h}=\alpha$
        by \auto{setext,dom_iff,diagonal_club_advance,pair_eq_iff}.
    Show that for every $\beta\in\dom{h}$ we have $h(\beta)\in\alpha$.
    \begin{subproof}
        Fix $\beta\in\dom{h}$.
        Hence $\beta\in\alpha$.
        Let $C=\indexedIntersection{\suc{\beta}}{F}{\alpha}$.
        $\ordinalUnbounded{C}{\alpha}$ by \auto.
        Hence $\clubUpper{C}{\beta}{\alpha}$ is inhabited
            by \obv{club_upper_inhabited}.
        Hence $\choice{\clubUpper{C}{\beta}{\alpha}}
            \in\clubUpper{C}{\beta}{\alpha}$ by \obv{choice_in}.
        Hence $\choice{\clubUpper{C}{\beta}{\alpha}}\in C$
            by \obv{club_upper}.
        $C\subseteq\alpha$ by \obv{ordinal_unbounded}.
        Hence $\choice{\clubUpper{C}{\beta}{\alpha}}\in\alpha$
            by \obv{subseteq}.
        $(\beta,\choice{\clubUpper{C}{\beta}{\alpha}})\in h$
            by \auto{diagonal_club_advance}.
        Hence $h(\beta)=\choice{\clubUpper{C}{\beta}{\alpha}}$
            by \obv{function_apply_intro}.
    \end{subproof}
    Hence $h$ is a function to $\alpha$.
    Follows by \auto{funs_intro}.
\end{proof}

\begin{proposition}\label{diagonal_club_advance_properties}
    Let $\alpha$ be an ordinal.
    Suppose for every $\rho\in\alpha$ we have
        $\ordinalUnbounded{
            \indexedIntersection{\suc{\rho}}{F}{\alpha}}{\alpha}$.
    Let $\beta\in\alpha$.
    Let $h=\diagonalClubAdvance{F}{\alpha}$.
    Then $\beta\in h(\beta)$ and for every $i\in\suc{\beta}$ we have
        $h(\beta)\in F(i)$.
\end{proposition}
\begin{proof}
    $h\in\funs{\alpha}{\alpha}$
        by \auto{diagonal_club_advance_type}.
    Hence $h$ is a function by \obv{funs_is_function}.
    Let $C=\indexedIntersection{\suc{\beta}}{F}{\alpha}$.
    $\ordinalUnbounded{C}{\alpha}$ by \auto.
    Hence $\clubUpper{C}{\beta}{\alpha}$ is inhabited
        by \obv{club_upper_inhabited}.
    Let $x=\choice{\clubUpper{C}{\beta}{\alpha}}$.
    $x\in\clubUpper{C}{\beta}{\alpha}$ by \auto{choice_in}.
    Hence $x\in C$ and $\beta\in x$ by \obv{club_upper}.
    $(\beta,x)\in h$ by \auto{diagonal_club_advance}.
    Hence $h(\beta)=x$ by \obv{function_apply_intro}.
    For every $i\in\suc{\beta}$ we have $x\in F(i)$
        by \auto{indexed_intersection_iff}.
    Follows by \auto.
\end{proof}

\begin{proposition}\label{diagonal_club_iteration_properties}
    Let $\alpha$ be an ordinal.
    Suppose for every $\rho\in\alpha$ we have
        $\ordinalUnbounded{
            \indexedIntersection{\suc{\rho}}{F}{\alpha}}{\alpha}$.
    Let $\beta\in\alpha$.
    Let $h=\diagonalClubAdvance{F}{\alpha}$.
    Let $a=h(\beta)$.
    Let $s=\iterationSet{h}{a}$.
    Then $s\in\funs{\naturals}{\alpha}$ and
        for every natural number $n$ we have
            $s(n)\in s(\suc{n})$ and
            for every $i\in\suc{s(n)}$ we have
                $s(\suc{n})\in F(i)$.
\end{proposition}
\begin{proof}
    $h\in\funs{\alpha}{\alpha}$
        by \auto{diagonal_club_advance_type}.
    $a\in\alpha$ by \auto{funs_type_apply}.
    $s\in\funs{\naturals}{\alpha}$
        by \auto{iteration_set_type}.
    Show that for every natural number $n$ we have
        $s(n)\in s(\suc{n})$ and
        for every $i\in\suc{s(n)}$ we have
            $s(\suc{n})\in F(i)$.
    \begin{subproof}
        Fix $n$.
        Assume $n$ is a natural number.
        $s(n)\in\alpha$ by \obv{funs_type_apply}.
        $s(\suc{n})=h(s(n))$
            by \auto{iteration_set_successor}.
        $s(n)\in h(s(n))$ and for every $i\in\suc{s(n)}$ we have
            $h(s(n))\in F(i)$
            by \obv{diagonal_club_advance_properties}.
        Follows by \auto.
    \end{subproof}
\end{proof}

\begin{proposition}\label{strict_ordinal_iteration_increasing}
    Let $\alpha$ be an ordinal.
    Let $h\in\funs{\alpha}{\alpha}$.
    Suppose for every $\rho\in\alpha$ we have $\rho\in h(\rho)$.
    Let $a\in\alpha$.
    Let $s=\iterationSet{h}{a}$.
    Let $m,n$ be natural numbers.
    Suppose $n\in m$.
    Then $s(n)\in s(m)$.
\end{proposition}
\begin{proof}
    $s\in\funs{\naturals}{\alpha}$
        by \auto{iteration_set_type}.
    Let $N=\{k\in\naturals\mid
        \text{for every $j\in k$ we have $s(j)\in s(k)$}\}$.
    $\emptyset\in N$ by \auto{omega_empty,emptyset}.
    Show that for every $k\in N$ we have $\suc{k}\in N$.
    \begin{subproof}
        Fix $k\in N$.
        Hence $k$ is a natural number and
            for every $j\in k$ we have $s(j)\in s(k)$.
        Hence $\suc{k}$ is a natural number by \obv{omega_successor}.
        $s(k)\in\alpha$ by \obv{funs_type_apply}.
        Hence $s(k)\in h(s(k))$.
        $s(\suc{k})=h(s(k))$
            by \auto{iteration_set_successor}.
        Hence $s(k)\in s(\suc{k})$.
        $s(\suc{k})\in\alpha$ by \obv{funs_type_apply}.
        Hence $s(\suc{k})$ is an ordinal
            by \obv{elem_of_ordinal_is_ordinal}.
        Show that for every $j\in\suc{k}$ we have
            $s(j)\in s(\suc{k})$.
        \begin{subproof}
            Fix $j\in\suc{k}$.
            Hence $j\in k$ or $j=k$ by \auto{suc_iff}.
            If $j=k$, then $s(j)\in s(\suc{k})$.
            If $j\in k$, then $s(j)\in s(k)$.
            Hence $s(j)\in s(\suc{k})$
                by \auto{ordinal_transitivity}.
        \end{subproof}
    \end{subproof}
    Hence $N$ is an inductive set by \obv{inductive_set}.
    Hence every natural number is an element of $N$
        by \obv{omega_minimal}.
    Hence $m\in N$.
\end{proof}

\begin{proposition}\label{strict_ordinal_iteration_limit}
    Let $\alpha$ be a limit ordinal.
    Suppose $\naturals\in\cofinality{\alpha}$.
    Let $h\in\funs{\alpha}{\alpha}$.
    Suppose for every $\rho\in\alpha$ we have $\rho\in h(\rho)$.
    Let $a\in\alpha$.
    Let $s=\iterationSet{h}{a}$.
    Let $x=\unions{\ran{s}}$.
    Then $x\in\alpha$ and $x$ is a limit ordinal and $a\in x$.
\end{proposition}
\begin{proof}
    $\alpha$ is an ordinal by \obv{limit_ordinal}.
    $s\in\funs{\naturals}{\alpha}$
        by \auto{iteration_set_type}.
    Hence $s$ is a function and $\dom{s}=\naturals$
        by \obv{funs_is_function,funs}.
    $x\in\alpha$ by \auto{small_domain_range_union_below}.
    $\ran{s}\subseteq\alpha$ by \obv{funs_ran}.
    Hence every element of $\ran{s}$ is an ordinal
        by \obv{subseteq,elem_of_ordinal_is_ordinal}.
    Hence $x$ is an ordinal by \auto{union_of_ordinals_is_ordinal}.
    $\emptyset$ is a natural number by \obv{omega_empty}.
    $s(\emptyset)=a$ by \auto{iteration_set_zero}.
    Hence $s(\emptyset)\in\alpha$.
    Hence $s(\emptyset)\in h(s(\emptyset))$.
    $s(\suc{\emptyset})=h(s(\emptyset))$
        by \auto{iteration_set_successor}.
    Hence $s(\emptyset)\in s(\suc{\emptyset})$.
    $\suc{\emptyset}$ is a natural number by \obv{omega_successor}.
    Hence $s(\suc{\emptyset})\in\ran{s}$ by \auto{fun_ran_iff}.
    Hence $s(\emptyset)\in x$ by \auto{unions_iff}.
    Hence $a\in x$.
    Hence $x\neq\emptyset$ by \auto{emptyset}.
    Hence $\emptyset\in x$
        by \obv{ordinal_empty_or_emptyset_elem}.
    Show that for every $\rho\in x$ we have $\suc{\rho}\in x$.
    \begin{subproof}
        Fix $\rho\in x$.
        Take $y\in\ran{s}$ such that $\rho\in y$
            by \auto{unions_iff}.
        Take $n\in\dom{s}$ such that $s(n)=y$
            by \auto{fun_ran_iff}.
        Hence $s(n)\in\alpha$ by \obv{funs_type_apply}.
        Hence $s(n)\in h(s(n))$.
        $s(\suc{n})=h(s(n))$
            by \auto{iteration_set_successor}.
        Hence $s(n)\in s(\suc{n})$.
        Hence $\rho$ is an ordinal by \obv{elem_of_ordinal_is_ordinal}.
        Hence $\suc{\rho}\subseteq s(n)$
            by \auto{ordinal_suc_subseteq}.
        $\suc{\rho}$ is an ordinal by \obv{suc_ordinal}.
        Hence $\suc{\rho}\in s(n)$ or $\suc{\rho}=s(n)$
            by \obv{ordinal_subseteq_implies_elem_or_eq}.
        $s(\suc{n})$ is an ordinal
            by \obv{funs_type_apply,omega_successor,elem_of_ordinal_is_ordinal}.
        Hence $\suc{\rho}\in s(\suc{n})$
            by \auto{ordinal_transitivity}.
        Hence $\suc{n}\in\dom{s}$ by \obv{funs,omega_successor}.
        Hence $s(\suc{n})\in\ran{s}$ by \auto{fun_ran_iff}.
        Follows by \auto{unions_iff}.
    \end{subproof}
    Hence $\emptyset\precedes x$.
    Hence $x$ is a limit ordinal by \obv{limit_ordinal}.
\end{proof}

\begin{proposition}\label{strict_ordinal_iteration_common_stage}
    Let $\alpha$ be an ordinal.
    Let $h\in\funs{\alpha}{\alpha}$.
    Suppose for every $\rho\in\alpha$ we have $\rho\in h(\rho)$.
    Let $a\in\alpha$.
    Let $s=\iterationSet{h}{a}$.
    Let $m,n$ be natural numbers.
    Then there exists a natural number $l$ such that
        $s(m)\in s(l)$ and $s(n)\in s(l)$.
\end{proposition}
\begin{proof}
    $m$ and $n$ are ordinals by \obv{natural_number_is_ordinal}.
    Let $k=m\union n$.
    $m\subseteq n$ or $n\subseteq m$
        by \obv{ordinal_subseteq_or_subseteq}.
    Hence $k=n$ or $k=m$
        by \auto{union_absorb_subseteq_left,union_absorb_subseteq_right}.
    Let $l=\suc{k}$.
    Hence $l$ is a natural number by \auto{omega_successor}.
    $m\subseteq k$ and $n\subseteq k$
        by \auto{union_upper_left,union_upper_right}.
    Hence $m\in k$ or $m=k$
        by \auto{ordinal_subseteq_implies_elem_or_eq}.
    Hence $m\in l$ by \auto{suc_iff}.
    Hence $s(m)\in s(l)$
        by \obv{strict_ordinal_iteration_increasing}.
    Hence $n\in k$ or $n=k$
        by \auto{ordinal_subseteq_implies_elem_or_eq}.
    Hence $n\in l$ by \auto{suc_iff}.
    Hence $s(n)\in s(l)$
        by \obv{strict_ordinal_iteration_increasing}.
    Follows by \auto.
\end{proof}

\begin{theorem}\label{diagonal_clubs_have_common_limit_above}
    Let $\alpha$ be a limit ordinal.
    Suppose $\naturals\in\cofinality{\alpha}$.
    Suppose for every $\rho\in\alpha$ we have
        $\ordinalUnbounded{
            \indexedIntersection{\suc{\rho}}{F}{\alpha}}{\alpha}$.
    Suppose for every $i\in\alpha$ we have
        $\ordinalClosed{F(i)}{\alpha}$.
    Let $\beta\in\alpha$.
    Then there exists $x\in\diagonalIntersection{F}{\alpha}$ such that
        $\beta\in x$ and $x$ is a limit ordinal.
\end{theorem}
\begin{proof}
    $\alpha$ is an ordinal by \obv{limit_ordinal}.
    Let $h=\diagonalClubAdvance{F}{\alpha}$.
    $h\in\funs{\alpha}{\alpha}$
        by \auto{diagonal_club_advance_type}.
    Show that for every $\rho\in\alpha$ we have $\rho\in h(\rho)$.
    \begin{subproof}
        Fix $\rho\in\alpha$.
        Follows by \auto{diagonal_club_advance_properties}.
    \end{subproof}
    Let $a=h(\beta)$.
    $a\in\alpha$ by \auto{funs_type_apply}.
    Let $s=\iterationSet{h}{a}$.
    $s\in\funs{\naturals}{\alpha}$
        by \auto{diagonal_club_iteration_properties}.
    Hence $s$ is a function and $\dom{s}=\naturals$
        by \obv{funs_is_function,funs}.
    Let $x=\unions{\ran{s}}$.
    $x\in\alpha$ and $x$ is a limit ordinal and $a\in x$
        by \obv{strict_ordinal_iteration_limit}.
    $\beta\in a$ by \obv{diagonal_club_advance_properties}.
    $x$ is an ordinal by \obv{limit_ordinal}.
    Hence $\beta\in x$ by \auto{ordinal_transitivity}.
    Show that for every $i\in x$ we have $x\in F(i)$.
    \begin{subproof}
        Fix $i\in x$.
        Hence $i\in\alpha$ by \auto{ordinal_transitivity}.
        $\ordinalClosed{F(i)}{\alpha}$.
        Show that $\ordinalUnbounded{F(i)\inter x}{x}$.
        \begin{subproof}
            $F(i)\inter x\subseteq x$ by \obv{subseteq,inter}.
            Show that for every $\theta\in x$ there exists
                $z\in F(i)\inter x$ such that $\theta\in z$.
            \begin{subproof}
                Fix $\theta\in x$.
                Take $y_m\in\ran{s}$ such that $\theta\in y_m$
                    by \auto{unions_iff}.
                Take $m\in\dom{s}$ such that $s(m)=y_m$
                    by \auto{fun_ran_iff}.
                Hence $m$ is a natural number.
                Take $y_n\in\ran{s}$ such that $i\in y_n$
                    by \auto{unions_iff}.
                Take $n\in\dom{s}$ such that $s(n)=y_n$
                    by \auto{fun_ran_iff}.
                Hence $n$ is a natural number.
                Take a natural number $l$ such that
                    $s(m)\in s(l)$ and $s(n)\in s(l)$
                    by \obv{strict_ordinal_iteration_common_stage}.
                $s(l)\in\alpha$ by \obv{funs_type_apply}.
                Hence $s(l)$ is an ordinal
                    by \obv{elem_of_ordinal_is_ordinal}.
                $\theta\in s(m)$ and $i\in s(n)$.
                Hence $\theta\in s(l)$ and $i\in s(l)$
                    by \auto{ordinal_transitivity}.
                Hence $i\in\suc{s(l)}$ by \auto{suc_iff}.
                Let $z=s(\suc{l})$.
                $z\in F(i)$
                    by \auto{diagonal_club_iteration_properties}.
                $s(l)\in z$
                    by \auto{diagonal_club_iteration_properties}.
                $z\in s(\suc{\suc{l}})$
                    by \auto{diagonal_club_iteration_properties,omega_successor}.
                $\suc{\suc{l}}$ is a natural number
                    by \obv{omega_successor}.
                Hence $\suc{\suc{l}}\in\dom{s}$ by \auto{funs}.
                Hence $s(\suc{\suc{l}})\in\ran{s}$
                    by \auto{fun_ran_iff}.
                Hence $z\in x$ by \auto{unions_iff}.
                Hence $z\in F(i)\inter x$ by \obv{inter_intro}.
                $z\in\alpha$ by \auto{funs_type_apply,omega_successor}.
                Hence $z$ is an ordinal
                    by \obv{elem_of_ordinal_is_ordinal}.
                Hence $\theta\in z$ by \auto{ordinal_transitivity}.
                Follows by \auto.
            \end{subproof}
            Follows by \obv{ordinal_unbounded}.
        \end{subproof}
        Follows by \obv{ordinal_closed}.
    \end{subproof}
    Hence $x\in\diagonalIntersection{F}{\alpha}$
        by \obv{diagonal_intersection}.
    Follows by \auto.
\end{proof}

\begin{theorem}\label{diagonal_club_intersection_from_prefixes}
    Let $\alpha$ be a limit ordinal.
    Suppose $\naturals\in\cofinality{\alpha}$.
    Suppose for every $\rho\in\alpha$ we have
        $\ordinalUnbounded{
            \indexedIntersection{\suc{\rho}}{F}{\alpha}}{\alpha}$.
    Suppose for every $i\in\alpha$ we have $\clubIn{F(i)}{\alpha}$.
    Then $\clubIn{\diagonalIntersection{F}{\alpha}}{\alpha}$.
\end{theorem}
\begin{proof}
    Let $D=\diagonalIntersection{F}{\alpha}$.
    $D\subseteq\alpha$ by \auto{diagonal_intersection,subseteq}.
    Show that $\ordinalUnbounded{D}{\alpha}$.
    \begin{subproof}
        Show that for every $\beta\in\alpha$ there exists $x\in D$
            such that $\beta\in x$.
        \begin{subproof}
            Fix $\beta\in\alpha$.
            For every $i\in\alpha$ we have
                $\ordinalClosed{F(i)}{\alpha}$
                by \auto{club_in_elim}.
            Take $x\in D$ such that $\beta\in x$ and
                $x$ is a limit ordinal
                by \auto{diagonal_clubs_have_common_limit_above}.
            Follows by \auto.
        \end{subproof}
        Follows by \obv{ordinal_unbounded}.
    \end{subproof}
    For every $i\in\alpha$ we have $\ordinalClosed{F(i)}{\alpha}$
        by \auto{club_in_elim}.
    $\ordinalClosed{\diagonalIntersection{F}{\alpha}}{\alpha}$
        by \obv{diagonal_intersection_closed}.
    Hence $\ordinalClosed{D}{\alpha}$.
    Follows by \auto{club_in}.
\end{proof}

\begin{theorem}\label{diagonal_club_intersection_regular}
    Let $\kappa$ be a regular cardinal.
    Suppose $\naturals\in\kappa$.
    Let $F\in\funs{\kappa}{\pow{\kappa}}$.
    Suppose for every $i\in\kappa$ we have $\clubIn{F(i)}{\kappa}$.
    Then $\clubIn{\diagonalIntersection{F}{\kappa}}{\kappa}$.
\end{theorem}
\begin{proof}
    $\kappa$ is an infinite cardinal and
        $\cofinality{\kappa}=\kappa$ by \obv{regular_cardinal}.
    Hence $\kappa$ is a limit ordinal
        by \obv{infinite_cardinal_is_limit_ordinal}.
    Hence $\kappa$ is an ordinal by \obv{limit_ordinal}.
    $\naturals\in\cofinality{\kappa}$.
    $F$ is a function and $\dom{F}=\kappa$
        by \obv{funs_is_function,funs}.
    Show that for every $\rho\in\kappa$ we have
        $\ordinalUnbounded{
            \indexedIntersection{\suc{\rho}}{F}{\kappa}}{\kappa}$.
    \begin{subproof}
        Fix $\rho\in\kappa$.
        Let $I=\suc{\rho}$.
        $I\in\kappa$ by \auto{suc_elem_limitordinal}.
        Hence $I\in\cofinality{\kappa}$.
        $\rho\in I$ by \obv{suc_intro_self}.
        Hence $I$ is inhabited by \auto.
        $\rho$ is an ordinal by \obv{elem_of_ordinal_is_ordinal}.
        Hence $I$ is an ordinal by \obv{suc_ordinal}.
        $I\subseteq\kappa$ by \obv{ordinal_elem_implies_subseteq}.
        Hence $I\subseteq\dom{F}$.
        For every $i\in I$ we have $\clubIn{F(i)}{\kappa}$
            by \obv{subseteq}.
        $\clubIn{\indexedIntersection{I}{F}{\kappa}}{\kappa}$
            by \obv{indexed_club_intersection}.
        Hence $\ordinalUnbounded{
            \indexedIntersection{I}{F}{\kappa}}{\kappa}$
            by \obv{club_in_elim}.
    \end{subproof}
    Follows by \obv{diagonal_club_intersection_from_prefixes}.
\end{proof}

\begin{proposition}\label{regressive_on_elim}
    Suppose $f$ is regressive on $S$.
    Then $f$ is a function and $\dom{f}=S$ and
        for every $\gamma\in S\setminus\{\emptyset\}$ we have
            $f(\gamma)\in\gamma$.
\end{proposition}
\begin{proof}
    Follows by \obv{regressive_on}.
\end{proof}

\begin{proposition}\label{fodor_avoiding_clubs_inhabited}
    Suppose $\nonstationaryIn{\fibre{f}{\eta}}{\kappa}$.
    Then $\fodorAvoidingClubs{f}{\eta}{\kappa}$ is inhabited.
\end{proposition}
\begin{proof}
    Take $C$ such that $\clubIn{C}{\kappa}$ and
        $\fibre{f}{\eta}\notmeets C$ by \obv{nonstationary_in}.
    $C\subseteq\kappa$ by \auto{club_in_elim}.
    Hence $C\in\pow{\kappa}$ by \obv{powerset_intro}.
    Hence $C\in\fodorAvoidingClubs{f}{\eta}{\kappa}$
        by \obv{fodor_avoiding_clubs}.
    Follows by \auto.
\end{proof}

\begin{theorem}\label{fodor_club_family_type}
    Suppose for every $\eta\in\kappa$ we have
        $\nonstationaryIn{\fibre{f}{\eta}}{\kappa}$.
    Let $D=\fodorClubFamily{f}{\kappa}$.
    Then $D\in\funs{\kappa}{\pow{\kappa}}$ and
        for every $\eta\in\kappa$ we have
            $\clubIn{D(\eta)}{\kappa}$ and
            $\fibre{f}{\eta}\notmeets D(\eta)$.
\end{theorem}
\begin{proof}
    $D$ is a relation by \auto{relation,fodor_club_family}.
    $D$ is right-unique
        by \auto{rightunique,fodor_club_family,pair_eq_iff}.
    $\dom{D}=\kappa$
        by \auto{setext,dom_iff,fodor_club_family,pair_eq_iff}.
    Show that for every $\eta\in\kappa$ we have
        $D(\eta)\in\pow{\kappa}$ and
        $\clubIn{D(\eta)}{\kappa}$ and
        $\fibre{f}{\eta}\notmeets D(\eta)$.
    \begin{subproof}
        Fix $\eta\in\kappa$.
        $\fodorAvoidingClubs{f}{\eta}{\kappa}$ is inhabited
            by \obv{fodor_avoiding_clubs_inhabited}.
        Let $C=\choice{\fodorAvoidingClubs{f}{\eta}{\kappa}}$.
        $C\in\fodorAvoidingClubs{f}{\eta}{\kappa}$
            by \auto{choice_in}.
        Hence $C\in\pow{\kappa}$ and $\clubIn{C}{\kappa}$ and
            $\fibre{f}{\eta}\notmeets C$
            by \obv{fodor_avoiding_clubs}.
        $(\eta,C)\in D$ by \auto{fodor_club_family}.
        Hence $D(\eta)=C$ by \obv{function_apply_intro}.
    \end{subproof}
    Hence $D$ is a function to $\pow{\kappa}$ by \auto.
    Hence $D\in\funs{\kappa}{\pow{\kappa}}$
        by \obv{funs_intro}.
    Follows by \auto.
\end{proof}

\begin{theorem}\label{fodor}
    Let $\kappa$ be a regular cardinal.
    Suppose $\naturals\in\kappa$.
    Suppose $S$ is stationary in $\kappa$.
    Let $f\in\funs{S}{\kappa}$.
    Suppose $f$ is regressive on $S$.
    Then there exists $\eta\in\kappa$ such that
        $\fibre{f}{\eta}$ is stationary in $\kappa$.
\end{theorem}
\begin{proof}
    $\kappa$ is an infinite cardinal and
        $\cofinality{\kappa}=\kappa$ by \obv{regular_cardinal}.
    Hence $\kappa$ is a limit ordinal
        by \obv{infinite_cardinal_is_limit_ordinal}.
    Hence $\kappa$ is an ordinal by \obv{limit_ordinal}.
    $\naturals\in\cofinality{\kappa}$.
    $S\subseteq\kappa$ by \obv{stationary_in_elim}.
    $f$ is a function and $\dom{f}=S$
        by \obv{funs_is_function,funs}.
    Suppose not.
    Show that for every $\eta\in\kappa$ we have
        $\nonstationaryIn{\fibre{f}{\eta}}{\kappa}$.
    \begin{subproof}
        Fix $\eta\in\kappa$.
        $\fibre{f}{\eta}\subseteq\dom{f}$
            by \obv{subseteq,function_fibre_iff}.
        Hence $\fibre{f}{\eta}\subseteq S$.
        Hence $\fibre{f}{\eta}\subseteq\kappa$
            by \obv{subseteq_transitive}.
        $\fibre{f}{\eta}$ is not stationary in $\kappa$ by \auto.
        Follows by \obv{nonstationary_iff_not_stationary}.
    \end{subproof}
    Let $D=\fodorClubFamily{f}{\kappa}$.
    $D\in\funs{\kappa}{\pow{\kappa}}$ and
        for every $\eta\in\kappa$ we have
            $\clubIn{D(\eta)}{\kappa}$ and
            $\fibre{f}{\eta}\notmeets D(\eta)$
        by \obv{fodor_club_family_type}.
    Let $E=\diagonalIntersection{D}{\kappa}$.
    $\clubIn{E}{\kappa}$
        by \auto{diagonal_club_intersection_regular}.
    Let $L=\limitOrdinalsBelow{\kappa}$.
    $\clubIn{L}{\kappa}$ by \auto{limit_ordinals_below_is_club}.
    $\clubIn{E\inter L}{\kappa}$
        by \obv{two_club_intersection}.
    Hence $S\meets(E\inter L)$ by \obv{stationary_in}.
    Take $\gamma$ such that $\gamma\in S$ and
        $\gamma\in E\inter L$ by \auto{disjoint}.
    Hence $\gamma\in E$ and $\gamma\in L$
        by \obv{inter_elim_left,inter_elim_right}.
    Hence $\gamma\in\kappa$ and $\gamma$ is a limit ordinal
        by \auto{limit_ordinals_below}.
    Hence $\emptyset\in\gamma$ by \obv{limit_ordinal}.
    Hence $\gamma\neq\emptyset$ by \auto{emptyset}.
    Hence $\gamma\notin\{\emptyset\}$ by \auto{singleton_iff}.
    Hence $\gamma\in S\setminus\{\emptyset\}$
        by \obv{setminus_intro}.
    Hence $f(\gamma)\in\gamma$ by \obv{regressive_on_elim}.
    Let $\eta=f(\gamma)$.
    Hence $\eta\in\kappa$ by \auto{ordinal_transitivity}.
    $\gamma\in D(\eta)$ by \auto{diagonal_intersection}.
    $\gamma\in\fibre{f}{\eta}$
        by \auto{function_fibre_iff}.
    $\fibre{f}{\eta}\notmeets D(\eta)$
        by \auto{fodor_club_family_type}.
    Contradiction by \auto{disjoint}.
\end{proof}

\begin{proposition}\label{fodor_values_inhabited}
    Let $\kappa$ be a regular cardinal.
    Suppose $\naturals\in\kappa$.
    Suppose $S$ is stationary in $\kappa$.
    Let $f\in\funs{S}{\kappa}$.
    Suppose $f$ is regressive on $S$.
    Then $\fodorValues{f}{\kappa}$ is inhabited.
\end{proposition}
\begin{proof}
    Take $\eta\in\kappa$ such that
        $\fibre{f}{\eta}$ is stationary in $\kappa$ by \obv{fodor}.
    Hence $\eta\in\fodorValues{f}{\kappa}$ by \obv{fodor_values}.
    Follows by \auto.
\end{proof}

\begin{proposition}\label{fodor_values_are_ordinals}
    Let $\kappa$ be an ordinal.
    For every $\eta\in\fodorValues{f}{\kappa}$ we have
        $\eta$ is an ordinal.
\end{proposition}
\begin{proof}
    Fix $\eta\in\fodorValues{f}{\kappa}$.
    Hence $\eta\in\kappa$ by \obv{fodor_values}.
    Follows by \obv{elem_of_ordinal_is_ordinal}.
\end{proof}

\begin{proposition}\label{fodor_value_mem}
    Let $\kappa$ be a regular cardinal.
    Suppose $\naturals\in\kappa$.
    Suppose $S$ is stationary in $\kappa$.
    Let $f\in\funs{S}{\kappa}$.
    Suppose $f$ is regressive on $S$.
    Then $\fodorValue{f}{\kappa}\in\fodorValues{f}{\kappa}$.
\end{proposition}
\begin{proof}
    $\kappa$ is a limit ordinal by \obv{stationary_in_elim}.
    Hence $\kappa$ is an ordinal by \obv{limit_ordinal}.
    $\fodorValues{f}{\kappa}$ is inhabited
        by \obv{fodor_values_inhabited}.
    Every element of $\fodorValues{f}{\kappa}$ is an ordinal
        by \obv{fodor_values_are_ordinals}.
    Follows by \auto{fodor_value,inters_of_ordinals_elem}.
\end{proof}

\begin{proposition}\label{fodor_value_lt}
    Let $\kappa$ be a regular cardinal.
    Suppose $\naturals\in\kappa$.
    Suppose $S$ is stationary in $\kappa$.
    Let $f\in\funs{S}{\kappa}$.
    Suppose $f$ is regressive on $S$.
    Then $\fodorValue{f}{\kappa}\in\kappa$.
\end{proposition}
\begin{proof}
    $\fodorValue{f}{\kappa}\in\fodorValues{f}{\kappa}$
        by \obv{fodor_value_mem}.
    Follows by \obv{fodor_values}.
\end{proof}

\begin{proposition}\label{fodor_fibre_eq}
    $\fodorFibre{f}{\kappa}
        =\fibre{f}{\fodorValue{f}{\kappa}}$.
\end{proposition}
\begin{proof}
    Follows by \obv{fodor_fibre}.
\end{proof}

\begin{proposition}\label{fodor_fibre_stationary}
    Let $\kappa$ be a regular cardinal.
    Suppose $\naturals\in\kappa$.
    Suppose $S$ is stationary in $\kappa$.
    Let $f\in\funs{S}{\kappa}$.
    Suppose $f$ is regressive on $S$.
    Then $\fodorFibre{f}{\kappa}$ is stationary in $\kappa$.
\end{proposition}
\begin{proof}
    $\fodorValue{f}{\kappa}\in\fodorValues{f}{\kappa}$
        by \obv{fodor_value_mem}.
    Follows by \auto{fodor_values,fodor_fibre_eq}.
\end{proof}

\begin{proposition}\label{fodor_fibre_subset_domain}
    Suppose $f$ is a function.
    Then $\fodorFibre{f}{\kappa}\subseteq\dom{f}$.
\end{proposition}
\begin{proof}
    Follows by \auto{fodor_fibre_eq,subseteq,function_fibre_iff}.
\end{proof}

\begin{proposition}\label{fodor_fibre_constant}
    Suppose $f$ is a function.
    For every $\beta\in\fodorFibre{f}{\kappa}$ we have
        $f(\beta)=\fodorValue{f}{\kappa}$.
\end{proposition}
\begin{proof}
    Fix $\beta\in\fodorFibre{f}{\kappa}$.
    Follows by \auto{fodor_fibre_eq,function_fibre_iff}.
\end{proof}

\begin{theorem}\label{fodor_value_and_fibre_properties}
    Let $\kappa$ be a regular cardinal.
    Suppose $\naturals\in\kappa$.
    Suppose $S$ is stationary in $\kappa$.
    Let $f\in\funs{S}{\kappa}$.
    Suppose $f$ is regressive on $S$.
    Then $\fodorValue{f}{\kappa}\in\kappa$ and
        $\fodorFibre{f}{\kappa}$ is stationary in $\kappa$ and
        $\fodorFibre{f}{\kappa}\subseteq S$ and
        for every $\beta\in\fodorFibre{f}{\kappa}$ we have
            $f(\beta)=\fodorValue{f}{\kappa}$.
\end{theorem}
\begin{proof}
    $f$ is a function and $\dom{f}=S$
        by \obv{funs_is_function,funs}.
    $\fodorValue{f}{\kappa}\in\kappa$ by \obv{fodor_value_lt}.
    $\fodorFibre{f}{\kappa}$ is stationary in $\kappa$
        by \obv{fodor_fibre_stationary}.
    $\fodorFibre{f}{\kappa}\subseteq\dom{f}$
        by \obv{fodor_fibre_subset_domain}.
    Hence $\fodorFibre{f}{\kappa}\subseteq S$.
    For every $\beta\in\fodorFibre{f}{\kappa}$ we have
        $f(\beta)=\fodorValue{f}{\kappa}$
        by \obv{fodor_fibre_constant}.
\end{proof}
